\chapter*{Abstract}
%\addcontentsline{toc}{chapter}{Abstract}
\label{chap:abstract}
\pagestyle{plain}
Coordinating heterogeneous teams of Unmanned Aerial Vehicles (UAVs) and Unmanned Ground Vehicles (UGVs) offers powerful operational synergies for complex missions including post-disaster search-and-rescue, environmental monitoring, and payload delivery. However, harnessing these complementary capabilities requires jointly solving two coupled NP-hard combinatorial sub-problems: task allocation and multi-vehicle routing under physical payload, battery capacity, capability compatibility, and mission deadline constraints. 

In this progress report, we formulate a bi-objective optimization model that minimizes a normalized trade-off between mission makespan and cumulative fleet energy consumption. To solve this problem at scale, a specialized Simulated Annealing (SA) metaheuristic is developed. The solver features an incremental cached route-state architecture providing a $5.25\times$ to $6.40\times$ speedup over full genome re-evaluation, a capability-aware greedy randomized construction heuristic, automatic temperature calibration, derived geometric cooling, and five domain-specific neighborhood operators (\texttt{relocate}, \texttt{swap}, \texttt{two\_opt}, \texttt{intra\_swap}, and \texttt{or\_opt}).

The algorithm is rigorously validated against an exact global optimization solver on small instances ($n=6\text{--}9$), achieving an exact 0.00\% optimality gap on $n=6,7,8$ and a 1.41\% mean gap on $n=9$ in under 1.6\,s. On a standard 8-customer Capacitated Vehicle Routing Problem (CVRP) benchmark, the solver consistently discovers the proven analytical global optimum (16.0\,m) across 100\% of trials. Comprehensive benchmark evaluations against Random Search, Greedy Construction, and First-Choice Hill Climbing demonstrate statistically decisive superiority ($p < 2 \times 10^{-6}$) and highlight the severe vulnerabilities of greedy heuristics to local traps.

Seven heterogeneous case studies ranging from 10 to 100 tasks and 2 to 10 vehicles are investigated. The solver achieves 100\% feasibility and stable convergence across all standard benchmarks, with empirical runtime scaling sub-linearly ($\mathcal{O}(n^{0.0985})$). In a high-stress endurance scenario with an 88\% reduction in UAV battery capacity (\texttt{battery75}), an extended-budget diagnostic confirms 100\% feasibility ($F=0.6049$) with zero constraint violations. Finally, systematic parameter sensitivity and ablation analyses establish the indispensable synergy between inter-vehicle workload rebalancing and intra-vehicle tour smoothing, while sweeping objective weights maps out the convex Pareto frontier between mission duration and fleet energy consumption.